% GENERATED FILE. Edit canonical lessons, then run npm run generate:books.
% canonical-source-hash: fnv1a64:8b1f18dfb2e78f1d
% canonical-lessons: ML-C253-almaram, ML-C253-plavu, ML-C253-rosappu, ML-C253-cemparatti, ML-C253-tulasi

\chapter{Banyan tree, Jackfruit tree, Rose, Hibiscus, Holy basil}
\label{ch:ml-banyan-tree-jackfruit-tree-rose-hibiscus-holy-basil}
\hlchaptermodality{253}

\begin{chapteropening}
\textbf{By the end of this chapter:} \emph{I can say a banyan tree, a jackfruit tree, a rose, a hibiscus and holy basil.}

Complete the last lesson of chapter 253: I can say a banyan tree, a jackfruit tree, a rose, a hibiscus and holy basil.
\end{chapteropening}

\section[\texorpdfstring{\ml{ആൽമരം} (ālmaraṁ)}{ആൽമരം (ālmaraṁ)}]{\ml{ആൽമരം} (ālmaraṁ) — a banyan tree}

\label{lesson:ML-C253-almaram}



\begin{quote}
Before the new one: say the Malayalam for a Vishu gift of money, then the Malayalam for a firecracker.
\end{quote}

\subsection*{You'll want to know: \ml{ആൽമരം}}
\textbf{\ml{ആൽമരം}} — \emph{ālmaraṁ} — \textquotedblleft{}a banyan tree\textquotedblright{}.

\begin{cousinweb}
\textbf{\ml{ആൽമരം}} (\emph{ālmaraṁ}) is native Dravidian: \emph{āl}, \textquotedblleft{}banyan\textquotedblright{} (DEDR 382, with Tamil \emph{āl} and Kannada \emph{āl}), and \emph{maram}, \textquotedblleft{}tree\textquotedblright{} (DEDR 4711).
\end{cousinweb}

\begin{grammarlens}[title={What you've built}]
One more word for this chapter.
\end{grammarlens}

\subsection*{Your turn}
\begin{itemize}
\raggedright
  \item \emph{Say it:} \emph{ālmaraṁ}
  \item \emph{Say it:} \emph{ālmaraṁ}, once more
  \item \emph{Say it:} say \emph{ālmaraṁ}
  \item \emph{Recall:} say the Malayalam for a Vishu gift of money, then the Malayalam for a firecracker, then say \emph{ālmaraṁ} again
\end{itemize}

\begin{tcolorbox}[breakable,colback=teal!4,colframe=teal!35!black,title={Before you move on}]
What does \ml{ആൽമരം} mean? (\textquotedblleft{}A banyan tree\textquotedblright{}.) Say it once more.
\end{tcolorbox}
\section[\texorpdfstring{\ml{പ്ലാവ്} (plāvŭ)}{പ്ലാവ് (plāvŭ)}]{\ml{പ്ലാവ്} (plāvŭ) — a jackfruit tree}

\label{lesson:ML-C253-plavu}



\begin{quote}
Before the new one: say the Malayalam for a firecracker, then the Malayalam for a banyan tree.
\end{quote}

\subsection*{You'll want to know: \ml{പ്ലാവ്}}
\textbf{\ml{പ്ലാവ്}} — \emph{plāvŭ} — \textquotedblleft{}a jackfruit tree\textquotedblright{}.

\begin{cousinweb}
\textbf{\ml{പ്ലാവ്}} (\emph{plāvŭ}) is native Dravidian, one word with Tamil \emph{palavu} (DEDR 3988). Kannada and Telugu use \emph{panasa}, the Sanskrit word for the same tree. DEDR 3988 compares Sanskrit \emph{panasa}.
\end{cousinweb}

\begin{grammarlens}[title={What you've built}]
One more word for this chapter.
\end{grammarlens}

\subsection*{Your turn}
\begin{itemize}
\raggedright
  \item \emph{Say it:} \emph{plāvŭ}
  \item \emph{Say it:} \emph{plāvŭ}, once more
  \item \emph{Say it:} say \emph{plāvŭ}
  \item \emph{Recall:} say the Malayalam for a firecracker, then the Malayalam for a banyan tree, then say \emph{plāvŭ} again
\end{itemize}

\begin{tcolorbox}[breakable,colback=teal!4,colframe=teal!35!black,title={Before you move on}]
What does \ml{പ്ലാവ്} mean? (\textquotedblleft{}A jackfruit tree\textquotedblright{}.) Say it once more.
\end{tcolorbox}
\section[\texorpdfstring{\ml{റോസാപ്പൂ} (ṟōsāppū)}{റോസാപ്പൂ (ṟōsāppū)}]{\ml{റോസാപ്പൂ} (ṟōsāppū) — a rose}

\label{lesson:ML-C253-rosappu}



\begin{quote}
Before the new one: say the Malayalam for a banyan tree, then the Malayalam for a jackfruit tree.
\end{quote}

\subsection*{You'll want to know: \ml{റോസാപ്പൂ}}
\textbf{\ml{റോസാപ്പൂ}} — \emph{ṟōsāppū} — \textquotedblleft{}a rose\textquotedblright{}.

\begin{cousinweb}
\textbf{\ml{റോസാപ്പൂ}} (\emph{ṟōsāppū}) joins \emph{rōsa}, from English \emph{rose} or Portuguese \emph{rosa}, to native \emph{pū}, \textquotedblleft{}flower\textquotedblright{} (DEDR 4345). DEDR 4345 compares Sanskrit \emph{puṣpa}, with a query.
\end{cousinweb}

\begin{grammarlens}[title={What you've built}]
One more word for this chapter.
\end{grammarlens}

\subsection*{Your turn}
\begin{itemize}
\raggedright
  \item \emph{Say it:} \emph{ṟōsāppū}
  \item \emph{Say it:} \emph{ṟōsāppū}, once more
  \item \emph{Say it:} say \emph{ṟōsāppū}
  \item \emph{Recall:} say the Malayalam for a banyan tree, then the Malayalam for a jackfruit tree, then say \emph{ṟōsāppū} again
\end{itemize}

\begin{tcolorbox}[breakable,colback=teal!4,colframe=teal!35!black,title={Before you move on}]
What does \ml{റോസാപ്പൂ} mean? (\textquotedblleft{}A rose\textquotedblright{}.) Say it once more.
\end{tcolorbox}
\section[\texorpdfstring{\ml{ചെമ്പരത്തി} (cemparatti)}{ചെമ്പരത്തി (cemparatti)}]{\ml{ചെമ്പരത്തി} (cemparatti) — a hibiscus}

\label{lesson:ML-C253-cemparatti}



\begin{quote}
Before the new one: say the Malayalam for a jackfruit tree, then the Malayalam for a rose.
\end{quote}

\subsection*{You'll want to know: \ml{ചെമ്പരത്തി}}
\textbf{\ml{ചെമ്പരത്തി}} — \emph{cemparatti} — \textquotedblleft{}a hibiscus\textquotedblright{}.

\begin{cousinweb}
\textbf{\ml{ചെമ്പരത്തി}} (\emph{cemparatti}) begins with native \emph{cem-}, \textquotedblleft{}red\textquotedblright{} (DEDR 1931, with Tamil \emph{ce-}); the origin of the rest is uncertain.
\end{cousinweb}

\begin{grammarlens}[title={What you've built}]
One more word for this chapter.
\end{grammarlens}

\subsection*{Your turn}
\begin{itemize}
\raggedright
  \item \emph{Say it:} \emph{cemparatti}
  \item \emph{Say it:} \emph{cemparatti}, once more
  \item \emph{Say it:} say \emph{cemparatti}
  \item \emph{Recall:} say the Malayalam for a jackfruit tree, then the Malayalam for a rose, then say \emph{cemparatti} again
\end{itemize}

\begin{tcolorbox}[breakable,colback=teal!4,colframe=teal!35!black,title={Before you move on}]
What does \ml{ചെമ്പരത്തി} mean? (\textquotedblleft{}A hibiscus\textquotedblright{}.) Say it once more.
\end{tcolorbox}
\section[\texorpdfstring{\ml{തുളസി} (tuḷasi)}{തുളസി (tuḷasi)}]{\ml{തുളസി} (tuḷasi) — holy basil, tulsi}

\label{lesson:ML-C253-tulasi}



\begin{quote}
Before the new one: say the Malayalam for a rose, then the Malayalam for a hibiscus.
\end{quote}

\subsection*{You'll want to know: \ml{തുളസി}}
\textbf{\ml{തുളസി}} — \emph{tuḷasi} — \textquotedblleft{}holy basil, tulsi\textquotedblright{}.

\begin{cousinweb}
\textbf{\ml{തുളസി}} (\emph{tuḷasi}) is the holy basil's name across the south: DEDR 3357 sets it beside Tamil \emph{tuḻāy} and Kannada \emph{toḷaci}, and compares Sanskrit \emph{tulasī}. Which way the word travelled is not settled.
\end{cousinweb}

\begin{grammarlens}[title={What you've built}]
One more word for this chapter.
\end{grammarlens}

\subsection*{Your turn}
\begin{itemize}
\raggedright
  \item \emph{Say it:} \emph{tuḷasi}
  \item \emph{Say it:} \emph{tuḷasi}, once more
  \item \emph{Say it:} say \emph{tuḷasi}
  \item \emph{Recall:} say the Malayalam for a rose, then the Malayalam for a hibiscus, then say \emph{tuḷasi} again
\end{itemize}

\begin{tcolorbox}[breakable,colback=teal!4,colframe=teal!35!black,title={Before you move on}]
What does \ml{തുളസി} mean? (\textquotedblleft{}Holy basil, tulsi\textquotedblright{}.) Say it once more.
\end{tcolorbox}
